\documentclass[usebib]{compactarticle}

\usepackage[english]{babel}

\title{Machine Learning for Automated Code Review}
\author{Jane Doe}
\subtitle{Master's Thesis Proposal}

\begin{document}

\section{Introduction to the Topic}

Software quality assurance relies heavily on code review processes.
Manual code review is time-consuming and prone to human error.
Machine learning approaches offer potential to automate aspects of code review and improve efficiency.

\section{Contribution to the State-of-the-Art}

Existing tools focus on style violations and syntactic issues.
This work extends automated code review to detect semantic bugs and architectural violations.
The proposed approach combines static analysis with deep learning models trained on historical code review data.

\section{Research Focus and Research Questions}

This thesis investigates how machine learning models can effectively identify code issues that require human attention in the review process.
The focus is on balancing accuracy with false-positive rates to create a practical tool for developers.

\begin{RQs}
\item How accurately can neural networks classify code changes into review-relevant categories?
\item Which code features are most predictive of issues that human reviewers identify?
\item How does model performance vary across different programming languages and project domains?
\end{RQs}

\section{Methodology for Research: Prototype Implementation}

\subsection{Previous Work}

The prototype builds on transformer-based models from the HuggingFace library \cite{Author24Generic}.
Git repositories are used to extract code changes and associated review comments.
The implementation leverages PyTorch for model training and scikit-learn for evaluation metrics.

\subsection{Requirements}

The prototype must process code diffs and output predictions with confidence scores.
It should support multiple programming languages, specifically Python and Java.
Performance requirements are not critical for the prototype; focus lies on prediction quality.
Real-time inference speed is not a requirement.

\subsection{Evaluation}

Evaluation uses a dataset of 10,000 reviewed pull requests from open-source projects.
Metrics include precision, recall, and F1-score for issue detection.
A case study compares model predictions against actual reviewer decisions.

\end{document}